% Part: proof-theory
% Chapter: normalization
% Section: introduction

\documentclass[../../../include/open-logic-section]{subfiles}

\begin{document}

\olfileid{pt}{nor}{int}

\olsection{Introduction}

If you have !!{prove}d a conditional~$!A \lif !B$, you might be in a
position to use it in !!a{proof} of~$!B$ by using~\Elim{\lif}. This
would of course also require !!a{proof}~$\delta_1$ of~$!A$. Your
!!{proof}~$\delta_1$ of~$!A \lif !B$, more likely than not, ends
in~\Intro{\lif}, which turns !!a{proof}~$\delta_2$ of~$!B$ from the
open assumption~$!A$ into one of~$!A \lif !B$. If it does, in your
final !!{proof} has the form
\[
\AxiomC{$\Discharge{!A}{x}$}
\RightLabel{$\delta_2$}
\DeduceC{$!B$}
\DischargeRule{\Intro{\lif}}{x}
\UnaryInfC{$!A \lif !B$}
\AxiomC{}
\RightLabel{$\delta_2$}
\DeduceC{$!A$}
\RightLabel{\Elim{\lif}}
\BinaryInfC{$!B$}
\DisplayProof
\]
Although your strategy is efficient in that you re-used something you
already had (!!a{proof} of~$!A \lif !B$), your !!{proof} isn't. It
seems like a detour to introduce $!A \lif !B$ only to immediately
eliminate it. There should be a more direct way to !!{prove}~$!B$.

In fact, there always is. ``Detours'' like the
above---!!a{introduction} rule immediately followed by
!!a{elimination} rule---can always be removed from natural
deduction~!!{proof}s. This is called \emph{normalization} and the
result is a \emph{normal} !!{proof} (or !!a{proof} \emph{in normal
form}).

The proof of this result is, like the cut-elimination theorem in the
sequent calculus, somewhat involved and requires us to consider lots
of cases. It is harder to prove for classical \Log{N1c} than it is for
intuitionistic \Log{N1i}. Just like cut-elimination shows that you
can't !!{prove} more using the \Cut{} rule in the sequent calculus,
normalization shows that you can't prove more things using detours
than you can while avoiding them. There are other similarities: normal
!!{proof} can be much longer than the shortest !!{proof} with detours
of the same result, just like !!{proof}s in the sequent calculus can
be much shorter with cuts than the shortest cut-free !!{proof}. The
sub-!!{formula} property, that you can always !!{prove} a result using
only sub-!!{formula}s of the result in your !!{proof}, follows from
normalization for natural deduction just as it follows from
cut-elimination for the sequent calculus.

\end{document}
